\begin{table}[t]
\centering
\small
\caption{Feature-group importance and redundancy diagnostics in the final scaling audit.}
\label{tab:importance-redundancy}
\begin{threeparttable}
\begin{adjustbox}{width=\textwidth}
\begin{tabular}{llcc}
\toprule
Diagnostic & Quantity & All seven & Competent subset \\
\midrule
Group permutation & Confidence/option drop & 0.148 & 0.163 \\
Group permutation & PCA-hidden drop & 0.104 & 0.120 \\
Redundancy diagnostic & Confidence/option AUROC & 0.770 & 0.799 \\
Redundancy diagnostic & PCA-hidden-only AUROC & 0.726 & 0.764 \\
Redundancy diagnostic & Combined AUROC & 0.755 & 0.786 \\
\bottomrule
\end{tabular}
\end{adjustbox}

\begin{tablenotes}[flushleft]\footnotesize\item \parbox{0.95\textwidth}{The confidence/option group causes the larger AUROC drop under permutation. In the fixed logistic redundancy diagnostic, adding PCA-hidden features can underperform because hidden features add variance without enough independent information.}\end{tablenotes}
\end{threeparttable}
\end{table}
